% ============================================================
% Section 5: Discussion
% ============================================================

\section{Discussion}
\label{sec:discussion}

\subsection{Why Multitask Outperforms cls\_only on Localization}

The most striking result in the ablation study is the collapse of Dice to near
zero when the segmentation head is removed (\textbf{cls\_only}: $0.002 \pm 0.008$
vs.\ $0.875 \pm 0.018$ for Multitask, $p = 1.3 \times 10^{-8}$). This is
expected by construction: without a segmentation head, no pixel-level predictions
are produced. However, the reciprocal comparison reveals a more nuanced
picture: \textbf{seg\_only} achieves a Dice score of $0.867 \pm 0.023$ that is
not significantly lower than Multitask ($p = 0.288$), yet its PR-AUC drops to
$0.728 \pm 0.053$ ($p = 1.5 \times 10^{-8}$). This asymmetry reveals that
\emph{it is the classification task that benefits most strongly from joint
training}: the shared encoder, regularized by the segmentation loss, learns
spatially-aware features that improve document-level detection far beyond what
a classification-only or threshold-on-mask approach achieves. Conversely,
the segmentation head performs almost equally well whether or not a dedicated
classification head is present, confirming that the U-Net decoder captures
sufficient forgery evidence on its own. The binary detection of an
identity document forgery requires aggregated global evidence that the
segmentation branch, trained to produce spatial maps, cannot reliably provide
without the explicit classification objective.

The comparison between Multitask and \textbf{unweighted\_losses} isolates the
contribution of HPO loss weighting. These two variants share the same architecture
and joint training objective; they differ only in whether $\lambda_{\text{mask}}$
is set by Pareto HPO (mean $2.39 \pm 0.43$ across folds) or fixed at $0.5$. The
Pareto-optimized weights yield a significantly higher Dice ($p = 1.3 \times 10^{-4}$,
$d = 1.01$, large effect) while maintaining equivalent PR-AUC ($p = 0.083$).
This finding confirms that the loss weighting is the principal driver of
localization quality and that a value near $\lambda_{\text{mask}} \approx 2.4$
is broadly optimal for this dataset.

\subsection{On the Pareto Advantage}

The scalar scalarization experiment (Section~\ref{subsec:scalar}) shows that
Pareto HPO does not outperform scalar selection in absolute metric terms on this
dataset ($n = 10$ folds, no significant differences). This result should be
interpreted carefully. The primary claim of the paper is not that Pareto HPO
produces higher numbers than scalar HPO, but that it provides a principled,
objective-agnostic selection mechanism that:

\begin{enumerate}
  \item \textbf{Eliminates the manual weight choice.} Scalar selection requires
  the practitioner to declare a single performance criterion (PR-AUC or Dice)
  as the optimization target. In a forensic deployment context, neither is
  clearly dominant: detection accuracy determines the security gate reliability,
  while localization accuracy determines the quality of forensic evidence.
  Pareto HPO avoids this choice.

  \item \textbf{Reveals the trade-off landscape.} The 500-point objective space
  in Figure~\ref{fig:pareto} provides a global picture of how PR-AUC and Dice
  interact across hyperparameter configurations. In particular, it shows that
  high-Dice configurations ($\geq 0.70$) consistently achieve PR-AUC $\geq 0.90$,
  confirming that the two objectives are broadly complementary (not conflicting)
  in this setting. This information is invisible to scalar HPO.

  \item \textbf{Is consistent with the evaluation protocol.} The DeepID 2025
  Challenge ranks systems independently on Track~1 and Track~2, assigning no
  implicit weighting. Minimizing the distance to the ideal point (Eq.~\ref{eq:ideal_point})
  mirrors this evaluation semantics.
\end{enumerate}

Numerically, the \textbf{significant difference against unweighted\_losses} in the
blind test ($d = 1.01$ on Dice) provides the empirical validation that HPO-based
weighting matters. The scalar vs.\ Pareto comparison is best understood as
demonstrating that, given a well-tuned $\lambda_{\text{mask}}$, the specific
selection criterion (distance to ideal point vs.\ maximizing a single metric)
has a second-order effect.

\subsection{Limitations and Future Work}

\paragraph{Single dataset.}
All results are obtained on FantasyID. While the dataset covers multiple document
types and two attack categories, it is a synthetic dataset generated at a single
research institute. Generalization to real-world identity documents, unseen
document designs, or attacks not present in FantasyID (e.g., physical tampering,
scanner artifacts) is unknown. Extension to the SIDTD
dataset~\cite{TODO_sidtd} and to the private PXL Vision test set would
substantially strengthen the external validity of the conclusions.

\paragraph{Architecture.}
The DocVerify encoder is a custom lightweight CNN trained from scratch. Transformer-based
encoders~\cite{TODO_vit} and models pre-trained on large-scale vision datasets
typically outperform scratch-trained CNNs when sufficient training data is
available. Revisiting the architecture with a stronger encoder is a natural
extension, particularly if cross-domain evaluation requires richer feature representations.

\paragraph{Challenge evaluation gap.}
The comparison in Table~\ref{tab:challenge} compares our system evaluated on an
internal holdout against challenge participants evaluated on the official test set.
Although the two sets are drawn from the same distribution (FantasyID), the results
are not directly comparable. An official submission to the challenge would
eliminate this caveat.

\paragraph{Variance on Track~2.}
The high variance in localization F1 ($\pm 0.096$) across seeds, driven by F1 on
bonafide images, represents a practical reliability concern. Future work should
investigate targeted training strategies (e.g., calibration of classification
confidence to reduce false-positive mask activations on bonafide images).
